\section{Introduction}
\label{sec:intro}

A medical VLM benchmark is usually published as a table. The table is the visible end of a pipeline: manifest construction, image preprocessing, prompt and schema design, batched inference, parsing, aggregation, and threshold fitting. When the pipeline is not reported, three things become invisible: the \emph{denominator} behind each number, the \emph{failure modes} that produced missing or constant rows, and the \emph{cost} that made the design feasible. In a study with 275,218 calls, an unreported denominator is not a presentation issue; it is a reproducibility failure.

This paper describes the engineering half of a three-layer chest X-ray VLM study. Its scientific results---four measurement reversals and a negative selective-trust result---are reported in a companion paper. We deliberately separate the two. The companion asks whether a reported robustness, reliance, or transfer number is identifiable. This paper asks whether a benchmark harness can make such a number \emph{reconstructible}: can a reader recover the call, the row, the denominator, and the exclusion rule that produced it, without a GPU and without re-running a model?

\textbf{What we built.} A harness with:
\begin{enumerate}[leftmargin=1.2em,itemsep=2pt,topsep=2pt]
\item \textbf{A per-layer design.} NIH runs a 19-cell synthetic grid, PadChest runs native shift plus a blank control, and Open-i runs prose generation plus a blank control. The budget is allocated to the question each layer can answer, not replicated uniformly (Section~\ref{sec:design}).
\item \textbf{A denominator census.} For any prediction file the census recomputes rows, unique keys, images, cells per image, and parse failures from the file alone, so a reported denominator can be checked by a third party without this harness (Section~\ref{sec:evaluation}).
\item \textbf{Provenance-bound resumability.} A run-level state file skips completed stages; a durable JSONL writer with an atomic progress sidecar resumes inside a stage; and the cross-modal and relay arms write a per-stage fingerprint of prompt, label set, manifest and runner, so a resumed stage whose configuration changed is refused instead of silently mixing two experiments. The main matrix predates the per-stage guard; its delivered state now carries a post-hoc recomputation of the run-level fingerprint, flagged as such (Sections~\ref{sec:execution} and~\ref{sec:evaluation}).
\item \textbf{A frozen artifact contract.} Manifests have required columns and exact label strings; predictions have a fixed JSON schema; manifests carry SHA-256 locks; a split-pure evaluation manifest exposes the scored 1,506 rows (Section~\ref{sec:artifacts}).
\item \textbf{A CPU-only audit toolkit.} Responsiveness, real-vs-real null, operating point, and contract coverage are computed from frozen predictions; no audit requires a model call (Section~\ref{sec:audit}).
\item \textbf{A measured comparison, not a self-report.} Five writers run the same 2,000-cell workload and nine faults are injected at the syscall boundary. We report what durability costs, what it buys, and which of five testable invariants each alternative satisfies (Section~\ref{sec:evaluation}).
\end{enumerate}

\textbf{Position against existing infrastructure.} General evaluation harnesses optimise for breadth and comparability, and their resumability is file-level: a result file that exists is treated as a finished unit. Workflow engines (Make, Snakemake, DVC) resume by content hash but know nothing about model calls or output contracts. Neither family publishes a fault-injection matrix, and neither asks whether the denominator of a reported number can be rebuilt from the released artifacts. We borrow from all three: file-level reuse from evaluation harnesses, content-hash identity from workflow engines, and write-ahead durability from databases. The contribution is the combination aimed at a specific failure --- an unreconstructible denominator --- together with the measurements that show where each mechanism actually helps.

\textbf{What this paper does not claim.} It does not propose a new VLM, a new benchmark score, or a new scientific finding, and it does not claim that durable writing is a novel technique. Append-and-fsync, write-ahead logs and content-addressed pipelines all predate this work. What is claimed is narrower and testable: that coverage can be audited from the artifact, that the durability guarantee has a measured price and a measured payoff against four alternatives, and that five invariants replacing ``zero error rows'' can actually fail.

\textbf{Why it matters.} Benchmark engineering failures are silent. A model that collapses to one answer still produces a metric. A resumed run with duplicate rows still produces a metric. A mis-encoded label column still produces a metric. The harness is designed around the principle that a contract violation must be \emph{loud}---an error row, a failed stage, a non-zero exit---rather than a plausible number.
